\chapter{Alcances del proyecto}

\section{Alcances}
El presente estudio aplicado al monitoreo de la calidad fisicoquímica del agua del Reservorio de Huayllarcocha, ubicado en el distrito de Cusco, provincia y departamento de Cusco, durante el periodo comprendido entre 04 de abril del año 2026 y 11 de agosto del año 2026.

En cuanto al muestreo, se contempla la identificación y selección de puntos de muestreo representativos dentro del reservorio, considerando su distribución espacial y accesibilidad. Asimismo, se ha diseñado un plan de monitoreo que establece una frecuencia de medición continua durante las 24 horas del día, a fin de capturar las variaciones temporales de los parámetros evaluados.

Respecto a las variables de estudio, la investigación se limita exclusivamente a la medición de pH y turbidez. El pH será determinado in situ mediante un pH-metro digital previamente calibrado, mientras que la turbidez se medirá con un turbidímetro, cuyos resultados se expresarán en unidades nefelométricas de turbidez (NTU).

Los datos registrados serán transmitidos en tiempo real a una red de monitoreo especialmente implementada para este fin. Dicha red contará con un dispositivo de adquisición de datos capaz de reconocer, procesar y almacenar automáticamente las señales provenientes de los equipos de medición, garantizando la integridad y trazabilidad de la información recolectada.

En la etapa de procesamiento y análisis, se realizará un tratamiento estadístico de los datos obtenidos, el cual incluirá medidas de tendencia central, dispersión y análisis de series temporales. Posteriormente, los valores registrados serán comparados con los Estándares de Calidad Ambiental (ECA) para agua establecidos en el Decreto Supremo Nº 004-2017-MINAM, así como con los parámetros de referencia de la Organización Mundial de la Salud (OMS). Esta comparación permitirá determinar si losniveles de pH y turbidez se encuentran dentro de los rangos considerados apropiados, según la categoría correspondiente de la normativa peruana.

Como productos del estudio, se contempla la elaboración de un informe técnico final que consigne los hallazgos, conclusiones y recomendaciones preliminares orientadas a posibles medidas de tratamiento. Dicho informe será presentado formalmente ante el equipo académico asesor.

\section{Limitaciones}
Análisis de parámetros microbiológicos, metales pesados u otros contaminantes diferentes al pH y turbidez.

Diseño ni ejecución de sistemas de tratamiento de agua.

Implementación física de acciones correctivas en las fuentes evaluadas.

Análisis de agua en zonas fuera del área geográfica delimitada.